\documentclass[11pt]{article}

\usepackage[T1]{fontenc}
\usepackage[utf8]{inputenc}
\usepackage{lmodern}
\usepackage[margin=1.08in]{geometry}
\usepackage{microtype}
\usepackage{amsmath,amssymb,amsthm,mathtools}
\usepackage{booktabs,tabularx,longtable,array}
\usepackage{enumitem}
\usepackage{xcolor}
\usepackage{listings}
\usepackage{fancyhdr}
\usepackage{hyperref}
\usepackage[nameinlink,noabbrev]{cleveref}

\hypersetup{
  colorlinks=true,
  linkcolor=blue!45!black,
  citecolor=green!35!black,
  urlcolor=blue!55!black,
  pdftitle={Alaoglu--Erdos: Congruence-Kernel Rigidity, Kummer Defect Primes, and Denominator Escape},
  pdfauthor={Research continuation, August 22, 2026}
}

\setlength{\parindent}{0pt}
\setlength{\parskip}{0.55em}
\setcounter{tocdepth}{2}
\setcounter{secnumdepth}{3}
\allowdisplaybreaks

\pagestyle{fancy}
\setlength{\headheight}{14pt}
\fancyhf{}
\fancyhead[L]{Alaoglu--Erd\H{o}s continuation}
\fancyhead[R]{Congruence kernels and defect primes}
\fancyfoot[C]{\thepage}
\renewcommand{\headrulewidth}{0.4pt}

\definecolor{codegray}{RGB}{246,246,246}
\definecolor{rulegray}{RGB}{120,120,120}
\definecolor{deepblue}{RGB}{30,60,120}
\definecolor{deepgreen}{RGB}{25,105,70}
\definecolor{deepred}{RGB}{145,45,45}

\lstdefinestyle{pythonstyle}{
  language=Python,
  basicstyle=\ttfamily\scriptsize,
  backgroundcolor=\color{codegray},
  frame=single,
  rulecolor=\color{rulegray},
  numbers=left,
  numberstyle=\tiny\color{rulegray},
  numbersep=7pt,
  showstringspaces=false,
  breaklines=true,
  columns=fullflexible,
  keepspaces=true,
  tabsize=4
}

\newtheorem{theorem}{Theorem}[section]
\newtheorem{proposition}[theorem]{Proposition}
\newtheorem{lemma}[theorem]{Lemma}
\newtheorem{corollary}[theorem]{Corollary}
\newtheorem{conjecture}[theorem]{Conjecture}
\newtheorem{question}[theorem]{Question}
\newtheorem{problem}[theorem]{Problem}
\theoremstyle{definition}
\newtheorem{definition}[theorem]{Definition}
\newtheorem{example}[theorem]{Example}
\newtheorem{remark}[theorem]{Remark}

\newcommand{\Z}{\mathbb Z}
\newcommand{\N}{\mathbb N}
\newcommand{\Q}{\mathbb Q}
\newcommand{\R}{\mathbb R}
\newcommand{\F}{\mathbb F}
\newcommand{\Gm}{\mathbb G_m}
\newcommand{\ord}{\operatorname{ord}}
\newcommand{\Gal}{\operatorname{Gal}}
\newcommand{\lcm}{\operatorname{lcm}}
\newcommand{\rad}{\operatorname{rad}}
\newcommand{\Span}{\operatorname{span}}
\newcommand{\Ker}{\operatorname{Ker}}
\newcommand{\im}{\operatorname{im}}
\newcommand{\vp}{v_p}
\newcommand{\thetaq}{\vartheta}
\newcommand{\Rel}{\mathcal K}
\newcommand{\Def}{\mathcal D}
\newcommand{\Sync}{\mathcal S}
\newcommand{\boxdelta}{\delta^{\square}}
\newcommand{\raydelta}{\delta^{\mathrm{ray}}}
\newcommand{\AEname}{Alaoglu--Erd\H{o}s}
\newcommand{\statusE}{\textcolor{deepgreen}{\textbf{[E]}}}
\newcommand{\statusC}{\textcolor{deepblue}{\textbf{[C]}}}
\newcommand{\statusD}{\textcolor{deepred}{\textbf{[D]}}}
\newcommand{\statusP}{\textcolor{orange!70!black}{\textbf{[P]}}}

\newenvironment{resultbox}[1]
{\begin{center}\begin{minipage}{0.94\textwidth}\hrule\vspace{0.45em}\textbf{#1}\par\vspace{0.25em}}
{\vspace{0.35em}\hrule\end{minipage}\end{center}}

\newenvironment{statusbox}
{\begin{center}\begin{minipage}{0.94\textwidth}\small\hrule\vspace{0.45em}}
{\vspace{0.35em}\hrule\end{minipage}\end{center}}

\begin{document}
\pagenumbering{gobble}

\begin{titlepage}
\centering
\vspace*{0.65cm}
{\Huge\bfseries The Alaoglu--Erd\H{o}s Conjecture\par}
\vspace{0.4cm}
{\LARGE\bfseries Congruence-Kernel Rigidity, Kummer Defect Primes,\par}
\vspace{0.14cm}
{\LARGE\bfseries and Denominator Escape\par}
\vspace{0.65cm}
{\large A nonduplicative continuation of the August 21, 2026 unified report\par}
\vspace{0.2cm}
{\large August 22, 2026\par}
\vfill
\begin{minipage}{0.9\textwidth}
\small
\textbf{Bottom line.}
The full conjecture is not proved here.  The new contribution is an exact local-global
framework that was not present in the attached report.  At a good prime $p$, the congruence
relations among $2$ and $3$ form a lattice
\[
 \Rel_p(2,3)=\{(r,s)\in\Z^2:2^r3^s\equiv1\pmod p\}.
\]
The inclusion $\Rel_p(2,3)\subseteq\Rel_p(M,A)$ is equivalent to the existence of a
\emph{synchronized local exponent} $e_p$ with
$M\equiv2^{e_p}$ and $A\equiv3^{e_p}\pmod p$.
Using the Corrales--Rodrig\'a\~nez--Schoof support theorem, almost-everywhere inclusion is
proved equivalent to the global diagonal form $(M,A)=(2^k,3^k)$.  Thus, for a fixed pair
$M=2^x$, $A=3^x$, integrality of $x$ is equivalent to synchronized local exponentiation at
all but finitely many primes.

The opposite direction is quantified.  If $B$ and $C$ are multiplicatively independent,
Kummer theory and Chebotarev produce a positive-density set of primes at which
$\ord_p(C)\nmid\ord_p(B)$.  For a suitable odd prime $\ell$, an explicit Frobenius class has
density $1/\ell^2$.  Consequently, a hypothetical nonintegral solution has positive-density
local defects on each of the three rays $(2,M)$, $(3,A)$, and $(6,MA)$.
Every such defect forces a prime divisor of a corresponding interpolation denominator; a
controlled squarefree part of that denominator has logarithm $N/\ell+o(N)$.  Finally, the
Bugeaud--Corvaja--Zannier gcd theorem shows that the exact two-point secant denominators on all
three rays have maximal exponential rate.

These results sharpen the remaining barrier rather than conceal it: the missing implication is
an archimedean-to-finite-place transfer strong enough to force synchronized local exponents,
whereas Kummer--Chebotarev shows that any nonintegral candidate would fail synchronization on
positive-density Frobenius sets.
\end{minipage}
\vfill
{\footnotesize Status labels: \statusE{} elementary proof in this report;
\statusC{} rigorous modulo a named classical theorem; \statusD{} exact finite computation;
\statusP{} proposed research target.\par}
\end{titlepage}

\pagenumbering{roman}
\begin{abstract}
Put $\thetaq=\log 3/\log 2$.  The \AEname{} conjecture asserts that
$2^x,3^x\in\Z$ implies $x\in\Z$.  The attached unified report already develops the
four-exponentials reduction, exact conditional semigroup structure, determinant and Smith-form
methods, $p$-adic transversality, sparse interpolation, special-function, spectral, moment,
information-theoretic, and canonical-product routes.  This continuation therefore does not
repeat those arguments.

For $p\nmid BC$, define the relation kernel
$\Rel_p(B,C)=\{(r,s)\in\Z^2:B^rC^s\equiv1\pmod p\}$.
We prove that
\[
 \Rel_p(2,3)\subseteq\Rel_p(M,A)
 \quad\Longleftrightarrow\quad
 \exists!\ e_p\pmod{|\langle2,3\rangle|}:
 (M,A)\equiv(2^{e_p},3^{e_p})\pmod p.
\]
Combining this finite cyclic-group lemma with the support theorem of
Corrales--Rodrig\'a\~nez and Schoof gives a global classification: for positive integers
$M,A$, almost-everywhere relation-kernel inclusion, almost-everywhere synchronized local
exponents, three elementary order-divisibility conditions, and the global identity
$(M,A)=(2^k,3^k)$ are equivalent.

We then prove a quantitative converse for multiplicatively independent pairs.  If the classes
of $B,C$ are independent modulo $\ell$th powers, the Kummer extension
$\Q(\zeta_\ell,B^{1/\ell},C^{1/\ell})$ contains a conjugacy-stable set of Frobenius elements
of relative size $1/\ell^2$ for which $B$ is an $\ell$th power modulo $p$ and $C$ is not.
These primes satisfy $\ord_p(C)\nmid\ord_p(B)$.  Hence every hypothetical nonintegral
\AEname{} pair has positive-density defect-prime sets.

A general collision-valuation lemma transfers those defects to interpolation denominators.
For the polynomial interpolating $B^j\mapsto C^j$ on $0\le j\le N$, its least integral
coefficient denominator is divisible by
\[
 \operatorname{lcm}_{1\le n\le N}
 \frac{B^n-1}{\gcd(B^n-1,C^n-1)}.
\]
The Kummer Frobenius set supplies a squarefree divisor with logarithm $N/\ell+o(N)$, and the
Bugeaud--Corvaja--Zannier theorem gives the exact secant-denominator rate
$\frac1n\log d_n(B,C)\to\log B$.  The report closes with a verified computation, a precise
failure analysis, and a Lean-oriented formalization plan.
\end{abstract}

\tableofcontents
\clearpage
\pagenumbering{arabic}

\section{Executive summary and claim ledger}
\label{sec:executive}

\subsection{The new exact statements}

\begin{resultbox}{Finite-prime synchronization criterion \statusE}
For every prime $p\nmid6MA$, relation-kernel inclusion
$\Rel_p(2,3)\subseteq\Rel_p(M,A)$ holds if and only if there is a unique residue class
$e_p$ modulo $h_p=|\langle2,3\rangle\subset\F_p^\times|$ such that
\[
 M\equiv2^{e_p}\pmod p,
 \qquad
 A\equiv3^{e_p}\pmod p.
\]
This is a finite cyclic-group theorem; no transcendence input is used.
\end{resultbox}

\begin{resultbox}{Global synchronized-power rigidity \statusC}
For positive integers $M,A$, the following are equivalent:
(i) $(M,A)=(2^k,3^k)$ for some $k\ge0$;
(ii) relation-kernel inclusion holds for all but finitely many primes;
(iii) synchronized local exponents exist for all but finitely many primes; and
(iv) for almost all primes,
\[
 \ord_p(M)\mid\ord_p(2),\quad
 \ord_p(A)\mid\ord_p(3),\quad
 \ord_p(MA)\mid\ord_p(6).
\]
The only non-elementary input is the 1997 support theorem of
Corrales--Rodrig\'a\~nez and Schoof.
\end{resultbox}

\begin{resultbox}{Positive-density Kummer defect primes \statusC}
If $B,C\in\Q_{>0}$ are multiplicatively independent, there is an odd prime $\ell$ and a set
$\Sigma_\ell(B,C)$ of rational primes of natural density $1/\ell^2$ such that
\[
 \ord_p(C)\nmid\ord_p(B)
 \qquad(p\in\Sigma_\ell(B,C)).
\]
The set is defined by the explicit residue-symbol conditions
\[
 p\equiv1\pmod\ell,
 \qquad B\in(\F_p^\times)^\ell,
 \qquad C\notin(\F_p^\times)^\ell.
\]
This is a direct Kummer--Chebotarev argument.
\end{resultbox}

\begin{resultbox}{Collision primes force denominator primes \statusE}
Let $Q\in\Q[X]$ interpolate integer data $Q(x_i)=y_i$ at distinct integer nodes, and let
$\delta$ be the least positive integer with $\delta Q\in\Z[X]$.  Then
\[
 v_p(\delta)\ge
 \max\{0,\,v_p(x_i-x_j)-v_p(y_i-y_j)\}
\]
for every pair $i\ne j$.  In particular, an input collision modulo $p$ not respected by the
outputs forces $p\mid\delta$.
\end{resultbox}

\begin{resultbox}{Ray-denominator and secant escape \statusE{} \statusC}
Let $\raydelta_N(B,C)$ be the least coefficient denominator of the degree-$N$ polynomial
interpolating $B^j\mapsto C^j$ for $0\le j\le N$.  Then
\[
 \operatorname{lcm}_{1\le n\le N}
 \frac{B^n-1}{\gcd(B^n-1,C^n-1)}
 \mid \raydelta_N(B,C).
\]
If $B,C$ are multiplicatively independent, the two-point denominator
$d_n(B,C)=(B^n-1)/\gcd(B^n-1,C^n-1)$ satisfies
\[
 \lim_{n\to\infty}\frac{\log d_n(B,C)}n=\log B.
\]
For the Kummer set above, the squarefree product of primes
$p\le\ell N+1$ in $\Sigma_\ell(B,C)$ divides $\raydelta_N(B,C)$ and has logarithm
$N/\ell+o(N)$.
\end{resultbox}

\subsection{What this does and does not accomplish}

The new route gives an exact local-global characterization of the desired conclusion and a
positive-density description of how every hypothetical counterexample must fail that
characterization.  It also converts the failure into explicit prime divisors of rational
interpolation denominators.  These are substantive additions: the attached report discusses
$p$-adic logarithmic rank, Fermat-quotient layers, and interpolation denominators as Smith
cokernel orders, but it does not use finite-field relation kernels, the support problem, or
Kummer Frobenius classes.

The final contradiction is still absent.  Denominator lower bounds alone do not help unless
one also constructs a rational or algebraic quantity whose archimedean height is smaller than
those denominators force.  The attached report repeatedly identifies exactly this denominator
or height tax.  The present approach supplies a more structured denominator divisor, not an
unjustified upper bound.

\subsection{Trust convention}

Throughout the report:
\begin{description}[leftmargin=1.2cm,style=nextline]
\item[\statusE] The statement is proved from elementary algebra, finite cyclic groups,
polynomial divisibility, or unique factorization in the text.
\item[\statusC] The deduction is complete modulo a named published theorem, principally the
support theorem, Chebotarev density, Kummer theory, or the Bugeaud--Corvaja--Zannier gcd bound.
\item[\statusD] Every finite decision is exact; floating-point values are used only to display
ratios of already exact counts and logarithmic sums.
\item[\statusP] The item is a research target or formalization plan, not a proved theorem.
\end{description}

\section{Nonduplication audit of the attached report}
\label{sec:audit}

The supplied source contains 20,483 lines and reconciles twenty-seven independent
continuations.  A thematic and keyword audit was performed before choosing the present route.
The point of this section is not to summarize that document again, but to make the scope
boundary explicit.

\begin{longtable}{p{0.28\textwidth}p{0.31\textwidth}p{0.33\textwidth}}
\toprule
Topic & Already in the attached report & Treatment here \\
\midrule
\endhead
Four exponentials and transcendence & Exact $2\times2$ reduction; rational and algebraic
exponent exclusions; the remaining four-exponentials wall. & Assumed as inherited context;
not reproved. \\
Conditional solution structure & Primitive common odd core, least generator, free monoid of
solutions and candidates, exact counts and block laws. & Used only to state that a
nonintegral candidate is multiplicatively independent on the three rays. \\
Archimedean determinants & Arbitrary-node repulsion, mixed-exponent semigroup determinants,
sparse interpolation, one-logarithm deficit. & No new archimedean determinant is claimed. \\
Finite-place methods & Iwasawa logarithms at $2$ and $3$, Fermat quotients, tropical layers,
structural-prime residuals, resultants, Smith profiles. & We work instead in the finite cyclic
group $\F_p^\times$ and track exact congruence-relation kernels. \\
Interpolation denominator & Identified as the order of a value vector in a Vandermonde
cokernel; structural-prime specializations and residual resultants. & We add a general
collision-valuation lower bound and identify positive-density Kummer primes that must divide
the denominator. \\
Gcd methods & The Bugeaud--Corvaja--Zannier theorem is mentioned as a generic obstruction to
manufacturing large common divisors. & We derive the exact asymptotic of a natural reduced
secant denominator and apply it on all three counterexample rays. \\
Kummer theory & Radical degree and entanglement literature are surveyed, with the conclusion
that they do not see the special real exponent. & We use a different Kummer object: a finite
extension detecting whether one rational number is an $\ell$th power modulo $p$ while another
is not. \\
Support problem & No occurrence of the Corrales--Rodrig\'a\~nez--Schoof theorem, congruence
relation kernels, or synchronized local exponents was found. & This is the central new branch. \\
\bottomrule
\end{longtable}

\begin{statusbox}
The phrase ``new'' in this report means new relative to the supplied unified report after this
audit.  Several ingredients are classical; the contribution is their exact assembly around
the \AEname{} pair, the positive-density defect theorem, and the denominator-transfer package.
\end{statusbox}

\section{External landscape audit as of August 22, 2026}
\label{sec:landscape}

The user's comparison with another model concerns rapidly changing public claims, so it was
checked separately from the mathematical argument.

\begin{table}[ht]
\centering
\small
\begin{tabularx}{\textwidth}{p{0.19\textwidth}X p{0.27\textwidth}}
\toprule
Claim & Publicly verifiable status & Evidence boundary \\
\midrule
Jacobian conjecture & A 2026 primary preprint by Gao gives counterexamples in every dimension
greater than two and describes Alp\"oge's earlier three-dimensional counterexample.  Thus the
general conjecture is false in dimensions $n\ge3$, while the two-dimensional case remains
open. & Primary arXiv sources are available; ``resolved'' must not be read as solving the
remaining $n=2$ case. \\
Hadamard matrices below $2000$ & Public trackers report exact constructions for the twelve
previously unresolved admissible orders.  Epoch AI marks order $668$ provisionally solved, and
OEIS records all twelve orders. & The constructions are finitely and exactly checkable, but a
conventional paper and full contribution account were not located in this audit. \\
Elliptic curve of rank at least $30$ & No public primary source, equation, or independence
certificate was located.  Epoch AI still listed the task as unsolved and the public record as
rank at least $29$. & A private result may exist, but it cannot be treated as verified without
the curve and thirty independent points. \\
New \AEname{} breakthrough & No public preprint, announcement, or Lean artifact matching the
reported private claim was located. & The private communication is not evidence available for
mathematical comparison. \\
\bottomrule
\end{tabularx}
\caption{Current-claim audit.  This table is contextual and is not used in any proof.}
\end{table}

The mathematical response to a competitive claim should therefore be a checkable theorem
package rather than a speculative declaration of victory.  The rest of this report is written
to that standard.

\section{Setup and elementary consequences of a hypothetical solution}
\label{sec:setup}

Let
\[
 \thetaq=\frac{\log3}{\log2}.
\]
Suppose
\begin{equation}
  M=2^x\in\Z_{>0},
  \qquad
  A=3^x\in\Z_{>0}.
  \label{eq:candidate}
\end{equation}
If $x<0$, then $0<2^x<1$, impossible for a positive integer.  Thus $x\ge0$.  If $x$ is
rational, writing $x=a/b$ in lowest terms and using unique factorization in
$2^a=M^b$ forces $b=1$.  Hence every nonintegral candidate is irrational; the attached report
proves the stronger statement that it is transcendental.

The following elementary independence observation will be used repeatedly.

\begin{lemma}[Three independent rays]
\label{lem:three-rays}
Assume \eqref{eq:candidate} and $x\notin\Z$.  Then each pair
\[
 (2,M),\qquad(3,A),\qquad(6,MA)
\]
is multiplicatively independent in $\Q_{>0}^{\times}$.
\end{lemma}

\begin{proof}
If $M^u=2^v$ for nonzero integers $u,v$, then
$2^{xu}=2^v$, so $x=v/u\in\Q$, hence $x\in\Z$, a contradiction.  The same proof applies to
$(3,A)$.  Finally, $(MA)^u=6^v$ gives $6^{xu}=6^v$, again forcing $x=v/u\in\Z$.
\end{proof}

\begin{remark}
The three rays have different logical roles.  Either of the first two is already sufficient
to recognize an integer exponent once an almost-everywhere order-divisibility statement is
known.  The third ray is needed for a symmetric classification of arbitrary pairs $(M,A)$:
it forces the two independently detected integer exponents to agree.
\end{remark}

\section{Finite-field relation kernels and synchronized exponents}
\label{sec:relkernels}

\subsection{The relation lattice}

Let $p$ be a prime not dividing $BC$.  Negative powers are then defined in
$\F_p^\times$.  Define
\begin{equation}
 \Rel_p(B,C)
 =\{(r,s)\in\Z^2:B^rC^s\equiv1\pmod p\}.
 \label{eq:relation-kernel}
\end{equation}
This is the kernel of the homomorphism
\[
 \Phi_{p;B,C}:\Z^2\longrightarrow\F_p^\times,
 \qquad(r,s)\longmapsto B^rC^s.
\]
It is a full-rank sublattice of $\Z^2$, with index
$|\langle B,C\rangle|$.

For the base pair, write
\[
 H_p=\langle2,3\rangle\subset\F_p^\times,
 \qquad h_p=|H_p|.
\]
Because $\F_p^\times$ is cyclic, $H_p$ is the unique subgroup of order $h_p$.

\subsection{Local factorization theorem}

\begin{theorem}[Relation-kernel factorization]
\label{thm:local-sync}
Let $p\nmid6MA$.  The following are equivalent.
\begin{enumerate}[label=(\roman*)]
\item $\Rel_p(2,3)\subseteq\Rel_p(M,A)$.
\item The rule $2^r3^s\mapsto M^rA^s$ is a well-defined group homomorphism
$H_p\to\F_p^\times$.
\item There exists an integer $e_p$, unique modulo $h_p$, such that
\begin{equation}
 M\equiv2^{e_p}\pmod p,
 \qquad
 A\equiv3^{e_p}\pmod p.
 \label{eq:sync-local}
\end{equation}
\end{enumerate}
\end{theorem}

\begin{proof}
The map $\Phi_{p;2,3}:\Z^2\to H_p$ is surjective, and its kernel is
$\Rel_p(2,3)$.  The map $\Phi_{p;M,A}:\Z^2\to\F_p^\times$ factors through the quotient
$\Z^2/\Rel_p(2,3)\cong H_p$ if and only if
$\Rel_p(2,3)\subseteq\Rel_p(M,A)$.  This proves (i)$\Leftrightarrow$(ii).

Let $g$ generate $\F_p^\times$, and write $H_p=\langle g^{(p-1)/h_p}\rangle$.  Every
homomorphism from the cyclic group $H_p$ to $\F_p^\times$ has image of order dividing
$h_p$.  Since a cyclic group has a unique subgroup of each order, that image lies in $H_p$.
Every endomorphism of $H_p$ is exponentiation by a unique class $e_p\pmod{h_p}$.  Evaluating
at the generators $2$ and $3$ gives \eqref{eq:sync-local}.  Conversely,
\eqref{eq:sync-local} defines the required exponentiation homomorphism.  Uniqueness follows
because $2$ and $3$ generate $H_p$.
\end{proof}

\begin{corollary}[Axis and diagonal shadows]
\label{cor:axis-shadows}
If the equivalent conditions of \cref{thm:local-sync} hold, then
\[
 \ord_p(M)\mid\ord_p(2),\qquad
 \ord_p(A)\mid\ord_p(3),\qquad
 \ord_p(MA)\mid\ord_p(6).
\]
\end{corollary}

\begin{proof}
The relation vectors $(\ord_p(2),0)$, $(0,\ord_p(3))$, and
$(\ord_p(6),\ord_p(6))$ lie in $\Rel_p(2,3)$.  Kernel inclusion sends them into
$\Rel_p(M,A)$, which is exactly the displayed divisibility.
\end{proof}

\begin{remark}
The converse of \cref{cor:axis-shadows} need not hold at a single prime.  Three selected
one-dimensional relations need not generate the full two-dimensional relation lattice.
Their power appears globally, where the support theorem turns almost-everywhere order
conditions into exact rational powers.
\end{remark}

\subsection{A local defect radius}

\begin{definition}
For $p\nmid6MA$, define the full local defect set
\[
 \Sync_p(M,A)=\Rel_p(2,3)\setminus\Rel_p(M,A).
\]
If it is nonempty, define its radius
\begin{equation}
 \lambda_p(M,A)=
 \min_{(r,s)\in\Sync_p(M,A)}\max\{|r|,|s|\}.
 \label{eq:defect-radius}
\end{equation}
If $\Sync_p(M,A)$ is empty, put $\lambda_p(M,A)=\infty$.
\end{definition}

By \cref{thm:local-sync}, $\lambda_p(M,A)=\infty$ exactly when a synchronized exponent
exists.  The quantity $\lambda_p$ measures how large an exponent box is needed to witness the
failure.

\section{Global support rigidity}
\label{sec:support}

\subsection{The classical support theorem}

We use the following theorem in its multiplicative-group form.

\begin{theorem}[Corrales--Rodrig\'a\~nez--Schoof]
\label{thm:CRS}
Let $u,v\in\Q^\times$.  Suppose that for all but finitely many rational primes $p$ and every
integer $n$,
\[
 u^n\equiv1\pmod p
 \quad\Longrightarrow\quad
 v^n\equiv1\pmod p.
\]
Then $v=u^k$ for some $k\in\Z$.
\end{theorem}

For primes not dividing $uv$, the hypothesis is equivalent to
$\ord_p(v)\mid\ord_p(u)$.  The theorem is a local-to-global detection result for membership in
a cyclic subgroup of $\Q^\times$.

\subsection{Three-ray classification}

\begin{theorem}[Global synchronized-power rigidity]
\label{thm:global-rigidity}
Let $M,A\in\Z_{>0}$.  The following are equivalent.
\begin{enumerate}[label=(\alph*)]
\item There exists $k\in\Z_{\ge0}$ such that
\begin{equation}
 M=2^k,
 \qquad A=3^k.
 \label{eq:global-diagonal}
\end{equation}
\item For all but finitely many primes $p$,
$\Rel_p(2,3)\subseteq\Rel_p(M,A)$.
\item For all but finitely many primes $p$, there is a synchronized exponent $e_p$ satisfying
\eqref{eq:sync-local}.
\item For all but finitely many primes $p$,
\begin{equation}
 \ord_p(M)\mid\ord_p(2),\qquad
 \ord_p(A)\mid\ord_p(3),\qquad
 \ord_p(MA)\mid\ord_p(6).
 \label{eq:three-support}
\end{equation}
\end{enumerate}
\end{theorem}

\begin{proof}
If \eqref{eq:global-diagonal} holds, then
$M^rA^s=(2^r3^s)^k$.  Thus every base relation is respected at every prime not dividing
$6MA$, proving (a)$\Rightarrow$(b).  The equivalence (b)$\Leftrightarrow$(c) is
\cref{thm:local-sync}, and (b)$\Rightarrow$(d) is \cref{cor:axis-shadows}.

Assume (d).  The first order divisibility and \cref{thm:CRS} give $M=2^a$ for some
$a\in\Z$.  Since $M$ is a positive integer and $2>1$, necessarily $a\ge0$.  Similarly,
$A=3^b$ with $b\ge0$.  The third order divisibility gives $MA=6^c$ with $c\ge0$.  Unique
factorization in
\[
 2^a3^b=2^c3^c
\]
forces $a=b=c$.  This is \eqref{eq:global-diagonal}.
\end{proof}

\begin{corollary}[Exact local-global reformulation of a fixed \AEname{} pair]
\label{cor:AE-localglobal}
Assume $M=2^x$ and $A=3^x$ are positive integers.  Then the following are equivalent:
\begin{enumerate}[label=(\roman*)]
\item $x\in\Z$;
\item synchronized local exponents exist for all but finitely many primes;
\item relation-kernel inclusion holds for all but finitely many primes.
\end{enumerate}
\end{corollary}

\begin{proof}
If $x\in\Z$, take $k=x$ in \cref{thm:global-rigidity}.  Conversely, either local condition
gives $(M,A)=(2^k,3^k)$ for some integer $k$.  Since the real exponential is injective,
$2^x=M=2^k$ implies $x=k$.
\end{proof}

\begin{corollary}[One-ray support criteria]
\label{cor:one-ray}
Under the hypotheses of \cref{cor:AE-localglobal}, each one of the following conditions alone
is equivalent to $x\in\Z$:
\[
 \ord_p(M)\mid\ord_p(2)\text{ for almost all }p,
\]
\[
 \ord_p(A)\mid\ord_p(3)\text{ for almost all }p,
\]
\[
 \ord_p(MA)\mid\ord_p(6)\text{ for almost all }p.
\]
\end{corollary}

\begin{proof}
The forward implication is immediate.  For the first condition, \cref{thm:CRS} gives
$M=2^k$, and then $2^x=2^k$ gives $x=k$.  The other two cases are identical.
\end{proof}

\begin{remark}
\Cref{cor:AE-localglobal} is an exact reformulation, not a proof of the conjecture.  Its value
is that it replaces an irrational real exponent by a family of finite cyclic-group
compatibility conditions and identifies precisely how much local compatibility would suffice.
\end{remark}

\section{Defect primes and their positive density}
\label{sec:defects}

\subsection{Order-defect sets}

\begin{definition}
For $B,C\in\Z_{>1}$, define
\begin{equation}
 \Def(B,C)=
 \{p:p\nmid BC,\ \ord_p(C)\nmid\ord_p(B)\}.
 \label{eq:defect-set}
\end{equation}
Equivalently, if $d=\ord_p(B)$, then $p\in\Def(B,C)$ exactly when
$C^d\not\equiv1\pmod p$.
\end{definition}

\begin{corollary}[Infinitely many defects]
\label{cor:infinite-defects}
If $B$ and $C$ are multiplicatively independent positive integers, then $\Def(B,C)$ is
infinite.
\end{corollary}

\begin{proof}
If the set were finite, then $\ord_p(C)\mid\ord_p(B)$ for almost all $p$.
\Cref{thm:CRS} would give $C=B^k$, contradicting multiplicative independence.
\end{proof}

By \cref{lem:three-rays}, a hypothetical nonintegral \AEname{} solution therefore has three
infinite defect sets
\begin{equation}
 \Def(2,M),\qquad \Def(3,A),\qquad \Def(6,MA).
 \label{eq:three-defect-sets}
\end{equation}
Each is contained in the set of primes at which synchronized exponentiation fails.

\subsection{Choosing a Kummer witness exponent}

The support theorem gives infinitude but not a quantitative set.  Kummer theory gives a
positive-density subset with a transparent Frobenius description.

Write a positive rational number as
$B=\prod_q q^{b_q}$ and $C=\prod_q q^{c_q}$, with finitely many nonzero integer exponents.
Multiplicative independence means that the two valuation vectors $(b_q)_q$ and $(c_q)_q$ are
linearly independent over $\Q$.  Hence there are primes $q_1,q_2$ with
\begin{equation}
 D=b_{q_1}c_{q_2}-b_{q_2}c_{q_1}\ne0.
 \label{eq:minor-D}
\end{equation}
Choose an odd prime $\ell\nmid D$.

\begin{lemma}[Independence modulo $\ell$th powers]
\label{lem:kummer-independence}
For an odd prime $\ell\nmid D$, the classes of $B$ and $C$ are linearly independent in
$\Q^\times/(\Q^\times)^\ell$.  They remain independent in
$F^\times/(F^\times)^\ell$, where $F=\Q(\zeta_\ell)$.
\end{lemma}

\begin{proof}
If $B^uC^v$ is an $\ell$th power in $\Q$ with $(u,v)\not\equiv(0,0)\pmod\ell$, then taking
$q_1$- and $q_2$-adic valuations gives a nonzero vector in the kernel modulo $\ell$ of the
matrix with determinant $D$, impossible because $\ell\nmid D$.

For the second statement, suppose a rational number $r$ is an $\ell$th power in
$F=\Q(\zeta_\ell)$.  If $r$ is not an $\ell$th power in $\Q$, Capelli's binomial
irreducibility criterion gives $[\Q(r^{1/\ell}):\Q]=\ell$.  But this field would be a
subfield of $F$, whose degree is $\ell-1$, impossible.  Thus the natural map
$\Q^\times/(\Q^\times)^\ell\to F^\times/(F^\times)^\ell$ is injective, and independence
persists.
\end{proof}

\subsection{The Kummer Frobenius set}

Let
\begin{equation}
 L=\Q(\zeta_\ell,B^{1/\ell},C^{1/\ell}).
 \label{eq:kummer-field}
\end{equation}
By \cref{lem:kummer-independence},
$[L:\Q(\zeta_\ell)]=\ell^2$, and
\[
 |\Gal(L/\Q)|=(\ell-1)\ell^2.
\]
Define $\Sigma_\ell(B,C)$ to be the set of unramified rational primes $p$ satisfying
\begin{equation}
 p\equiv1\pmod\ell,
 \qquad
 B\in(\F_p^\times)^\ell,
 \qquad
 C\notin(\F_p^\times)^\ell.
 \label{eq:kummer-set}
\end{equation}

\begin{theorem}[Kummer defect-density theorem]
\label{thm:kummer-density}
With $B,C,\ell$ as above, the set $\Sigma_\ell(B,C)$ has natural density
\begin{equation}
 \delta(\Sigma_\ell(B,C))=\frac1{\ell^2}.
 \label{eq:density-one-over-l2}
\end{equation}
Every prime in this set lies in $\Def(B,C)$.  Consequently,
\begin{equation}
 \underline\delta\bigl(\Def(B,C)\bigr)\ge\frac1{\ell^2}.
 \label{eq:defect-density-lower}
\end{equation}
\end{theorem}

\begin{proof}
Over $F=\Q(\zeta_\ell)$, Kummer theory identifies
$\Gal(L/F)$ with $(\Z/\ell\Z)^2$.  Choose coordinates so that
$(a,b)$ sends
\[
 B^{1/\ell}\longmapsto\zeta_\ell^aB^{1/\ell},
 \qquad
 C^{1/\ell}\longmapsto\zeta_\ell^bC^{1/\ell}.
\]
Inside $\Gal(L/\Q)$, the subset
\[
 \mathcal C=\{(0,b):b\in(\Z/\ell\Z)^\times\}
\]
is stable under conjugacy: the cyclotomic Galois group multiplies both Kummer coordinates by
the same nonzero scalar.  It has $\ell-1$ elements.  Chebotarev's density theorem therefore
gives density
\[
 \frac{|\mathcal C|}{|\Gal(L/\Q)|}
 =\frac{\ell-1}{(\ell-1)\ell^2}
 =\frac1{\ell^2}.
\]
The associated splitting conditions are exactly \eqref{eq:kummer-set}: Frobenius is trivial on
$F$, fixes $B^{1/\ell}$, and acts nontrivially on $C^{1/\ell}$.

Now write $p-1=\ell^t m$ with $\ell\nmid m$.  In the cyclic group $\F_p^\times$, an
$\ell$th power has multiplicative order whose $\ell$-adic valuation is at most $t-1$.
A non-$\ell$th power has order with $\ell$-adic valuation exactly $t$.  Hence
\[
 v_\ell(\ord_p(B))\le t-1<t=v_\ell(\ord_p(C)),
\]
so $\ord_p(C)\nmid\ord_p(B)$.  This proves the containment and density lower bound.
\end{proof}

\begin{corollary}[Explicit witness for a non-power integer]
\label{cor:explicit-witness}
Suppose $M$ is a positive integer that is not a power of $2$.  Let $q$ be any odd prime
factor of $M$, put $t=v_q(M)$, and choose any odd prime $\ell\nmid t$.  Then
$\Sigma_\ell(2,M)$ has density $1/\ell^2$ and is contained in $\Def(2,M)$.
\end{corollary}

\begin{proof}
The $2\times2$ valuation minor using the primes $2$ and $q$ is $t$, so
\cref{lem:kummer-independence,thm:kummer-density} apply.
\end{proof}

\begin{corollary}[Positive-density failure for a hypothetical counterexample]
\label{cor:candidate-density}
If \eqref{eq:candidate} holds with $x\notin\Z$, then each of the three sets in
\eqref{eq:three-defect-sets} has positive lower density.  In particular, synchronized local
exponentiation fails on a positive-density set of primes.
\end{corollary}

\begin{proof}
Apply \cref{thm:kummer-density} to each multiplicatively independent pair in
\cref{lem:three-rays}.  An axis order defect prevents relation-kernel inclusion by
\cref{cor:axis-shadows}, hence prevents a synchronized exponent by \cref{thm:local-sync}.
\end{proof}

\subsection{Weighted prime asymptotics}

Chebotarev's prime number theorem gives the stronger weighted form
\begin{equation}
 \sum_{\substack{p\le X\\p\in\Sigma_\ell(B,C)}}\log p
 =\frac{X}{\ell^2}+o(X).
 \label{eq:chebotarev-theta}
\end{equation}
This will turn the positive-density Frobenius set into an explicit squarefree divisor of an
interpolation denominator.

\section{Interpolation denominators detect local relation defects}
\label{sec:denominators}

\subsection{A general collision-valuation lemma}

Let $x_0,\ldots,x_N$ be distinct integers and $y_0,\ldots,y_N$ integers.  Let
$Q\in\Q[X]$ be the unique polynomial of degree at most $N$ with $Q(x_i)=y_i$, and define
\begin{equation}
 \delta(Q)=\min\{d\in\Z_{>0}:dQ\in\Z[X]\}.
 \label{eq:min-denom}
\end{equation}

\begin{theorem}[Collision-valuation bound]
\label{thm:collision-valuation}
For every prime $p$ and every $i\ne j$ with $y_i\ne y_j$,
\begin{equation}
 v_p(\delta(Q))
 \ge
 \max\{0,\,v_p(x_i-x_j)-v_p(y_i-y_j)\}.
 \label{eq:collision-valuation}
\end{equation}
In particular,
\begin{equation}
 x_i\equiv x_j\pmod p,
 \quad y_i\not\equiv y_j\pmod p
 \quad\Longrightarrow\quad
 p\mid\delta(Q).
 \label{eq:collision-prime}
\end{equation}
\end{theorem}

\begin{proof}
Put $P=\delta(Q)Q\in\Z[X]$.  The polynomial $P(X)-P(x_j)$ is divisible by $X-x_j$ in
$\Z[X]$.  Evaluating at $x_i$ gives
\[
 x_i-x_j\mid P(x_i)-P(x_j)
 =\delta(Q)(y_i-y_j).
\]
Taking $p$-adic valuations yields
\[
 v_p(x_i-x_j)
 \le v_p(\delta(Q))+v_p(y_i-y_j),
\]
which is \eqref{eq:collision-valuation}.  The final implication is the case where the
left-hand difference has valuation at least one and the right-hand output difference has
valuation zero.
\end{proof}

\begin{remark}
The attached report identifies the interpolation denominator as an order in a Vandermonde
cokernel.  \Cref{thm:collision-valuation} gives a direct local witness for that order: two
rows of the evaluation matrix coalesce modulo $p^a$, while their prescribed values coalesce
to smaller depth.  No Smith-form computation is needed to certify the resulting denominator
factor.
\end{remark}

\subsection{Geometric-ray interpolation}

For integers $B,C>1$, let $Q_N^{B,C}$ be the unique degree-$N$ polynomial satisfying
\begin{equation}
 Q_N^{B,C}(B^j)=C^j
 \qquad(0\le j\le N),
 \label{eq:ray-interpolation}
\end{equation}
and put
\[
 \raydelta_N(B,C)=\delta(Q_N^{B,C}).
\]
Define the reduced secant denominator
\begin{equation}
 d_n(B,C)=\frac{B^n-1}{\gcd(B^n-1,C^n-1)}.
 \label{eq:secant-denom}
\end{equation}

\begin{theorem}[Secant divisors of the full ray denominator]
\label{thm:secant-divides}
For every $1\le n\le N$,
\begin{equation}
 d_n(B,C)\mid\raydelta_N(B,C).
 \label{eq:dn-divides}
\end{equation}
Consequently,
\begin{equation}
 \lcm_{1\le n\le N}d_n(B,C)
 \mid\raydelta_N(B,C).
 \label{eq:lcm-divides}
\end{equation}
\end{theorem}

\begin{proof}
Apply \cref{thm:collision-valuation} to the interpolation points
$(x_0,y_0)=(1,1)$ and $(x_n,y_n)=(B^n,C^n)$.  For every prime $p$,
\[
 v_p(\raydelta_N(B,C))
 \ge\max\{0,v_p(B^n-1)-v_p(C^n-1)\}
 =v_p(d_n(B,C)).
\]
This proves \eqref{eq:dn-divides}; taking the primewise maximum over $n$ gives
\eqref{eq:lcm-divides}.
\end{proof}

\begin{corollary}[Order defects divide the ray denominator]
\label{cor:defect-divides-ray}
If $p\in\Def(B,C)$ and $\ord_p(B)\le N$, then
$p\mid\raydelta_N(B,C)$.
\end{corollary}

\begin{proof}
Set $n=\ord_p(B)$.  Then $p\mid B^n-1$ but $p\nmid C^n-1$, so $p\mid d_n(B,C)$.
Apply \cref{thm:secant-divides}.
\end{proof}

\subsection{A controlled squarefree denominator divisor}

For a Kummer witness prime $\ell$, define
\begin{equation}
 R_N(B,C;\ell)=
 \prod_{\substack{p\le\ell N+1\\p\in\Sigma_\ell(B,C)}}p.
 \label{eq:kummer-divisor}
\end{equation}

\begin{theorem}[Kummer squarefree divisor]
\label{thm:kummer-denominator}
For every $N\ge1$,
\begin{equation}
 R_N(B,C;\ell)\mid\raydelta_N(B,C).
 \label{eq:kummer-divides-denom}
\end{equation}
Moreover,
\begin{equation}
 \log R_N(B,C;\ell)=\frac{N}{\ell}+o(N).
 \label{eq:kummer-divisor-asymp}
\end{equation}
\end{theorem}

\begin{proof}
For $p\in\Sigma_\ell(B,C)$, the element $B$ is an $\ell$th power modulo $p$.  Hence
\[
 \ord_p(B)\mid\frac{p-1}{\ell}.
\]
If $p\le\ell N+1$, then $\ord_p(B)\le N$.  By
\cref{thm:kummer-density}, $p\in\Def(B,C)$, so
\cref{cor:defect-divides-ray} gives $p\mid\raydelta_N(B,C)$.  The primes in the product are
distinct, proving \eqref{eq:kummer-divides-denom}.

Finally, use \eqref{eq:chebotarev-theta} with $X=\ell N+1$:
\[
 \log R_N(B,C;\ell)
 =\sum_{\substack{p\le\ell N+1\\p\in\Sigma_\ell(B,C)}}\log p
 =\frac{\ell N+1}{\ell^2}+o(N)
 =\frac N\ell+o(N).
\]
\end{proof}

\begin{remark}
The divisor in \cref{thm:kummer-denominator} is not merely large; every prime in it belongs
to one specified Frobenius class and witnesses a one-unit loss in $\ell$-adic order.  This
structure is the feature that could interact with a future finite-place determinant or
resultant construction.  Its exponential size by itself is not enough for a contradiction.
\end{remark}

\subsection{The two-dimensional candidate grid}

Let
\[
 \Omega_N=\{0,1,\ldots,N\}^2,
 \qquad
 x_{a,b}=2^a3^b,
 \qquad
 y_{a,b}=M^aA^b.
\]
The nodes $x_{a,b}$ are distinct.  Let $Q_N^\square\in\Q[X]$ be the unique polynomial of
degree below $(N+1)^2$ with
$Q_N^\square(x_{a,b})=y_{a,b}$, and let
$\boxdelta_N(M,A)$ be its least coefficient denominator.

\begin{theorem}[Box denominator detects relation defects]
\label{thm:box-detects}
Let $p\nmid6MA$.
\begin{enumerate}[label=(\roman*)]
\item If $\lambda_p(M,A)\le N$, then $p\mid\boxdelta_N(M,A)$.
\item If $p\nmid\boxdelta_{p-2}(M,A)$, then
$\Rel_p(2,3)\subseteq\Rel_p(M,A)$ and a synchronized local exponent exists.
\item Equivalently, failure of synchronized local exponentiation at $p$ forces
\begin{equation}
 p\mid\boxdelta_{p-2}(M,A).
 \label{eq:p-divides-box}
\end{equation}
\end{enumerate}
\end{theorem}

\begin{proof}
For (i), choose $(r,s)\in\Sync_p(M,A)$ with
$\max(|r|,|s|)\le N$.  Write
$r=r_+-r_-$ and $s=s_+-s_-$ with nonnegative parts.  The two grid points
$(r_+,s_+)$ and $(r_-,s_-)$ lie in $\Omega_N$.  Their input nodes are congruent modulo $p$
because $2^r3^s\equiv1$, but their output values are not because
$M^rA^s\not\equiv1$.  Apply \cref{thm:collision-valuation}.

For (ii), take any relation $(r,s)\in\Rel_p(2,3)$ and reduce each coordinate modulo $p-1$ to
$r',s'\in\{0,\ldots,p-2\}$.  Then
$x_{r',s'}\equiv1=x_{0,0}\pmod p$.  If $p$ does not divide the box denominator, the collision
must be respected, so $M^{r'}A^{s'}\equiv1$.  Reduction modulo $p-1$ does not change powers in
$\F_p^\times$, hence $(r,s)\in\Rel_p(M,A)$.  Apply \cref{thm:local-sync}.
Statement (iii) is the contrapositive.
\end{proof}

\begin{remark}
The scale $p-2$ is exact but too large for a direct global determinant argument.  A major
quantitative target is to prove that many defect primes have
$\lambda_p(M,A)\ll(\log p)^C$ or even $p^{1/2-\eta}$.  Kummer ray defects only give the axis
bound $\lambda_p\le(p-1)/\ell$.  Improving this relation-vector scale is one of the clearest
new bottlenecks exposed by the kernel language.
\end{remark}

\section{Maximal secant-denominator escape from the gcd theorem}
\label{sec:BCZ}

The two-point interpolant through $(1,1)$ and $(B^n,C^n)$ is
\[
 L_n(X)=1+\frac{C^n-1}{B^n-1}(X-1).
\]
Its least coefficient denominator is exactly $d_n(B,C)$ from
\eqref{eq:secant-denom}: both coefficients have common denominator
$(B^n-1)/\gcd(B^n-1,C^n-1)$ after reduction.

\begin{theorem}[Bugeaud--Corvaja--Zannier]
\label{thm:BCZ}
If $B,C>1$ are multiplicatively independent integers, then for every $\varepsilon>0$ and all
sufficiently large $n$,
\begin{equation}
 \gcd(B^n-1,C^n-1)\le e^{\varepsilon n}.
 \label{eq:BCZ}
\end{equation}
\end{theorem}

\begin{corollary}[Exact secant-denominator rate]
\label{cor:secant-rate}
If $B,C>1$ are multiplicatively independent, then
\begin{equation}
 \lim_{n\to\infty}\frac1n\log d_n(B,C)=\log B.
 \label{eq:secant-rate}
\end{equation}
\end{corollary}

\begin{proof}
By definition,
\[
 \log d_n(B,C)
 =\log(B^n-1)-\log\gcd(B^n-1,C^n-1).
\]
The first term divided by $n$ tends to $\log B$.  By \cref{thm:BCZ}, the second term divided
by $n$ tends to zero.  This proves the limit.
\end{proof}

\begin{corollary}[Three-ray maximal escape for a counterexample]
\label{cor:three-ray-escape}
If \eqref{eq:candidate} holds with $x\notin\Z$, then
\begin{align}
 \lim_{n\to\infty}\frac1n\log d_n(2,M)&=\log2,\\
 \lim_{n\to\infty}\frac1n\log d_n(3,A)&=\log3,\\
 \lim_{n\to\infty}\frac1n\log d_n(6,MA)&=\log6.
 \label{eq:three-rates}
\end{align}
In particular, the reduced secant denominators escape at the largest exponential rate allowed
by their input denominators on all three rays.
\end{corollary}

\begin{proof}
Use \cref{lem:three-rays,cor:secant-rate}.
\end{proof}

\begin{remark}
This is stronger than the qualitative statement that generic common divisors are too small.
It identifies an exact rational object naturally attached to the candidate interpolation
problem and computes its exponential rate.  It also explains why a proof based on unexpectedly
small secant denominators cannot work: a hypothetical counterexample forces the opposite,
namely asymptotically maximal reduced denominators.
\end{remark}

\section{Consequences specialized to the \AEname{} problem}
\label{sec:specialization}

Collecting the preceding results gives a compact conditional profile.

\begin{theorem}[Local arithmetic profile of a hypothetical nonintegral solution]
\label{thm:profile}
Assume $M=2^x$ and $A=3^x$ are integers with $x\notin\Z$.  Then:
\begin{enumerate}[label=(\arabic*)]
\item \statusE{} The pairs $(2,M)$, $(3,A)$, and $(6,MA)$ are multiplicatively independent.
\item \statusC{} Each ray has a positive-density set of order-defect primes; for suitable odd
primes $\ell_2,\ell_3,\ell_6$, the corresponding explicit Kummer subsets have densities
$1/\ell_2^2$, $1/\ell_3^2$, and $1/\ell_6^2$.
\item \statusC{} Synchronized local exponentiation fails on a positive-density set of primes.
\item \statusE{} Every ray defect $p$ with $\ord_p(B)\le N$ divides the ray interpolation
denominator $\raydelta_N(B,C)$.
\item \statusC{} Each ray denominator contains a controlled squarefree Frobenius divisor with
logarithm $N/\ell_i+o(N)$.
\item \statusC{} Each two-point secant denominator has its maximal exponential rate, as in
\eqref{eq:three-rates}.
\item \statusE{} At every prime of full synchronization failure,
$p\mid\boxdelta_{p-2}(M,A)$.
\end{enumerate}
\end{theorem}

\begin{proof}
Items (1)--(7) are respectively
\cref{lem:three-rays,thm:kummer-density,cor:candidate-density,cor:defect-divides-ray,thm:kummer-denominator,cor:three-ray-escape,thm:box-detects}.
\end{proof}

\subsection{An exact dichotomy}

For each candidate pair, there are only two possibilities.

\begin{enumerate}[label=\textbf{Case \Alph*.},leftmargin=2.2cm]
\item $x\in\Z$.  Then one global exponent $k=x$ works at every good prime; every relation is
respected; the ray interpolants eventually become monomials $X^k$ and have denominator one.
\item $x\notin\Z$.  Then synchronized local exponentiation fails on positive-density Kummer
sets; the ray and box interpolation denominators acquire the exact defect divisors described
above; and the two-point denominators have maximal exponential rate.
\end{enumerate}

This dichotomy is unusually sharp, but it is not self-contradictory.  The second case is
arithmetically coherent as far as the present theorems go.

\section{Exact computational reconnaissance}
\label{sec:computation}

The computation in this section is illustrative, not evidentiary for the infinite theorems.
Its purposes are to test the finite-group equivalences, display the role of the third ray, and
check convergence toward the Kummer density.

\subsection{Synthetic pair tests}

For a good prime $p$, the code computes multiplicative orders exactly, checks for an exponent
$e_p$ modulo $h_p$, and finds the first order defect on each ray.  Integer-power pairs pass by
theorem; mismatched powers show why the $(6,MA)$ ray is needed.

\begin{table}[ht]
\centering
\scriptsize
\setlength{\tabcolsep}{3.4pt}
\begin{tabular}{@{}cccccc@{}}
\toprule
$(M,A)$
& \shortstack{first\\$\Def(2,M)$}
& \shortstack{first\\$\Def(3,A)$}
& \shortstack{first\\$\Def(6,MA)$}
& \shortstack{first full\\failure}
& \shortstack{synchronized / good\\$p\le10^4$} \\
\midrule
$(4,9)$   & none & none & none & none & $1227/1227$ \\
$(8,27)$  & none & none & none & none & $1227/1227$ \\
$(4,27)$  & none & none & $5$ & $5$ & $1/1227$ \\
$(8,9)$   & none & none & $5$ & $5$ & $1/1227$ \\
$(5,13)$  & $7$ & $11$ & $7$ & $7$ & $1/1225$ \\
$(10,38)$ & $7$ & $13$ & $7$ & $7$ & $2/1225$ \\
\bottomrule
\end{tabular}
\caption{Exact modular reconnaissance.  ``None'' for the first two rows is a theorem, not
merely a search statement.  For display consistency, the search cutoff for a first defect was
$10^5$.}
\end{table}

The pairs $(4,27)$ and $(8,9)$ separately satisfy the first two ray identities globally:
$4=2^2$, $27=3^3$ and $8=2^3$, $9=3^2$.  The diagonal ray detects the mismatch immediately,
because $MA$ is not a common power of $6$.

\subsection{A Kummer density experiment}

For $(B,C,\ell)=(2,5,3)$, the Kummer set is
\[
 p\equiv1\pmod3,
 \qquad 2^{(p-1)/3}\equiv1\pmod p,
 \qquad 5^{(p-1)/3}\not\equiv1\pmod p.
\]
Its predicted density is $1/9=0.111111\ldots$.

\begin{table}[ht]
\centering
\small
\begin{tabular}{rrrrr}
\toprule
$X$ & $\#\Sigma_3(2,5)$ & good primes & count ratio &
$X^{-1}\sum_{p\le X,\,p\in\Sigma}\log p$ \\
\midrule
$10^3$ & $17$ & $166$ & $0.102409638554$ & $0.095640842040$ \\
$10^4$ & $139$ & $1227$ & $0.113284433578$ & $0.112447409348$ \\
$10^5$ & $1062$ & $9590$ & $0.110740354536$ & $0.110329538804$ \\
$10^6$ & $8701$ & $78496$ & $0.110846412556$ & $0.110666164775$ \\
\bottomrule
\end{tabular}
\caption{Exact prime and residue decisions; decimal ratios are display-only.  Both columns
approach the Chebotarev prediction $1/9$.}
\end{table}

The exact script has SHA-256 digest
\[
 \texttt{bf14ae4af00e825968c3bddde08ca9ee073a9b12c734fb88ac8c3c4e49215f33}.
\]
Its captured output has digest
\[
 \texttt{aedb75029c41695873d3a2203ca0a5b38d31997c885a186ddb0ecc66a09ab31b}.
\]
The complete source appears in \cref{app:code}.

\section{Why this still stops short}
\label{sec:barrier}

\subsection{The missing real-to-local bridge}

The new local-global theorem says that it would be enough to prove, from the real equalities
$M=2^x$ and $A=3^x$, that almost every congruence relation among $2$ and $3$ is respected by
$M$ and $A$.  There is no known mechanism for doing this.  The identity
\[
 \log M\,\log3=\log A\,\log2
\]
is archimedean.  Reducing it modulo $p$ is meaningless, and replacing real logarithms by
$p$-adic logarithms changes the kernel.  The attached report's Iwasawa-logarithm analysis
makes this mismatch explicit at the structural primes.

Kummer--Chebotarev makes the obstacle sharper: if a nonintegral candidate exists, the desired
local synchronization does not merely fail at sporadic primes.  It fails on positive-density,
explicit Frobenius sets.  A successful proof must therefore force an almost-everywhere local
property that is strongly incompatible with ordinary multiplicative independence.

\subsection{Why denominator divisibility is not yet a contradiction}

The collision lemma proves that defect primes divide interpolation denominators.  But a large
denominator is compatible with a rational polynomial of large height.  The full geometric and
square-grid interpolation matrices already have very large Vandermonde denominators.  The new
Kummer divisor is a structured part of that denominator, not an upper bound for the remaining
part and not a small algebraic integer.

A product-formula contradiction would require one of the following additional inputs:
\begin{enumerate}[label=(\alph*)]
\item a primitive numerator or resultant whose archimedean height is
$o(\exp(N/\ell))$ while the Kummer divisor survives;
\item cancellation that removes the generic node denominator but preserves the defect-prime
part;
\item a low-degree auxiliary polynomial whose value class cannot absorb the positive-density
Frobenius divisor; or
\item a relation-radius theorem that detects many defect primes in an exponent box much
smaller than $p$.
\end{enumerate}
None is presently proved.

\subsection{A geometry-of-numbers attempt and its scale failure}

One might try to choose a short relation $(r,s)\in\Rel_p(2,3)$ for which
$|r\log2+s\log3|$ is very small.  A lattice rectangle argument can trade coordinate size
against the real linear form, but only polynomially in the lattice index.  To turn
\[
 M^rA^s=(2^r3^s)^x
\]
into an equality of rationals by archimedean smallness, one needs an error exponentially small
in $|r|+|s|$, because clearing negative exponents introduces denominators of size roughly
$M^{|r|}A^{|s|}$.  Polynomial lattice approximation cannot reach that scale.  This recovers,
in relation-kernel language, the shrinking-target and normalization-tax barriers already
found by other routes in the attached report.

\subsection{Why positive density alone is insufficient}

Positive density gives
$\sum_{p\le X,p\in\Sigma}\log p\asymp X$, but a determinant built from all exponents up to
$N$ can have archimedean height of order $\exp(cN^2)$ or worse.  The Kummer product captured by
$N$ has logarithm only $\asymp N$.  To win, one must either compress the archimedean object to
linear logarithmic height or amplify each Frobenius defect beyond one squarefree prime factor.
This is a quantitative formulation of the next barrier.

\section{A prioritized research program}
\label{sec:program}

\subsection{Target I: short defect relations}

\begin{problem}[Short-defect problem \statusP]
For a hypothetical nonintegral candidate $(M,A)$, prove that there are infinitely many, or a
positive-density set of, primes $p$ for which
\[
 \lambda_p(M,A)\le p^{1/2-\eta}
\]
for some fixed $\eta>0$.  A polylogarithmic bound would be substantially stronger.
\end{problem}

A theorem of this kind would move defect-prime divisibility into much smaller square-grid
interpolants, where the archimedean height might be competitive.  It asks for more than a
short relation in $\Rel_p(2,3)$: the relation must lie outside $\Rel_p(M,A)$.

\subsection{Target II: primitive denominator extraction}

\begin{problem}[Frobenius-stable primitive denominator \statusP]
Construct from the data
$2^a3^b\mapsto M^aA^b$ an algebraic integer or primitive resultant $R_N$ such that:
\begin{enumerate}[label=(\roman*)]
\item every Kummer defect prime $p\le cN$ divides $R_N$;
\item the generic Vandermonde denominator cancels from $R_N$; and
\item $\log|R_N|=o(N)$ at the distinguished archimedean embedding.
\end{enumerate}
\end{problem}

The attached report's boundary-dilation resultants are the closest existing objects, but their
known heights dominate their local divisibility.  The Kummer condition suggests selecting
residuals by power-residue symbols rather than only by the structural primes $2$ and $3$.

\subsection{Target III: character-filtered interpolation}

For a fixed $\ell$, the indicator that $B$ is an $\ell$th power and $C$ is not can be written
using order-$\ell$ multiplicative characters at primes $p\equiv1\pmod\ell$.  This suggests a
character-filtered determinant or resultant whose local initial form vanishes on the
$B$-residue component but not on the $C$ component.  The exact Frobenius density gives the
expected supply of primes.  The unresolved point is global coherence: characters live in
different residue fields and cannot simply be added across primes.

\subsection{Target IV: support theorem with quantitative certificates}

The qualitative support theorem says that if $C$ is not a power of $B$, some prime detects the
failure.  The Kummer proof in \cref{sec:defects} provides a positive-density certificate after
choosing $\ell$.  An effective Chebotarev theorem, conditionally on GRH or unconditionally with
weaker bounds, can bound the least such prime in terms of the discriminant of
$\Q(\zeta_\ell,B^{1/\ell},C^{1/\ell})$.  For the \AEname{} problem this does not bound the unknown
integers $M,A$, but it can make every finite search branch fully explicit once a size range is
fixed.

\subsection{Target V: combine the three rays in one Kummer extension}

A joint Kummer extension generated by roots of
$2,3,6,M,A,MA$ may contain Frobenius classes that impose simultaneous defects or prescribed
patterns across the three rays.  Dependencies among these rational numbers determine the
actual Kummer rank.  Computing the exact conjugacy-stable class sizes could yield stronger
multi-ray denominator divisors than the union of three independent one-ray arguments.  This is
a finite linear-algebra problem over $\F_\ell$ once the prime-exponent vectors of $M,A$ are
known, but those vectors are unknown in the conjecture.

\section{Lean formalization blueprint}
\label{sec:lean}

The rival claim reportedly emphasizes Lean.  The present route naturally separates into a
large elementary finite component and three classical black-box theorems.

\subsection{Elementary modules}

The following units should be formalizable directly with existing finite-group and polynomial
infrastructure.

\begin{enumerate}[label=\texttt{K\arabic*.},leftmargin=1.45cm]
\item \texttt{FiniteRelationKernel}: define the homomorphism
$\Z\times\Z\to\F_p^\times$ and its kernel.
\item \texttt{SynchronizedExponent}: prove \cref{thm:local-sync} using cyclicity of
$\F_p^\times$ and the universal property of quotient groups.
\item \texttt{AxisDiagonalSupport}: derive the three order-divisibility shadows.
\item \texttt{CollisionDenominator}: formalize polynomial difference divisibility and
\cref{thm:collision-valuation} in a discrete valuation setting.
\item \texttt{GeometricRayDenominator}: prove
$d_n(B,C)\mid\raydelta_N(B,C)$ without constructing Lagrange coefficients.
\item \texttt{BoxDefect}: formalize the positive/negative-part translation of a relation
vector and \cref{thm:box-detects}.
\item \texttt{ExactRecon}: optionally verify the finite tables by reflection or native
decision procedures.
\end{enumerate}

A schematic theorem signature for the central finite lemma is:

\begin{lstlisting}[basicstyle=\ttfamily\footnotesize,frame=single,backgroundcolor=\color{codegray}]
theorem relKernel_le_iff_exists_syncExponent
    (p : Nat) [Fact p.Prime]
    (M A : ZMod p) (hgood : 2 * 3 * M * A != 0) :
    relKernel (2 : ZMod p) 3 <= relKernel M A <->
      Exists e : ZMod (orderOfSubgroup (2 : ZMod p) 3),
        M = 2 ^ e.val /\ A = 3 ^ e.val
\end{lstlisting}

The exact encoding will depend on whether exponent vectors are represented in $\Z^2$, a free
abelian group, or a finitely supported function type.

\subsection{Classical theorem interfaces}

Three inputs should be isolated as explicit axiomatized interfaces until independently
formalized:

\begin{enumerate}[label=\texttt{C\arabic*.},leftmargin=1.45cm]
\item the Corrales--Rodrig\'a\~nez--Schoof support theorem for $\Q^\times$;
\item Kummer degree and Chebotarev density for the extension in \eqref{eq:kummer-field};
\item the Bugeaud--Corvaja--Zannier subexponential gcd theorem.
\end{enumerate}

Everything after these interfaces is short.  In particular, the global three-ray theorem is
just three applications of the support interface plus unique factorization, and the secant
rate is one line after the gcd interface.

\subsection{Formalization order}

The most efficient order is K4, K5, K1, K2, K3, K6, followed by the global wrappers.  This
produces useful denominator lemmas even before the deep number-theoretic interfaces are
available.  Every theorem should retain the status distinction between a kernel-checked
finite statement and a deduction from a classical theorem imported as a hypothesis.

\section{Conclusions}
\label{sec:conclusion}

This continuation does not claim a proof of the \AEname{} conjecture.  It does produce a new,
coherent theorem package that is both mathematically exact and sharply connected to the final
obstruction.

The main structural advance is the congruence-kernel reformulation:
\[
 x\in\Z
 \quad\Longleftrightarrow\quad
 \text{synchronized local exponents exist for almost every }p
\]
for every fixed integral-output pair $(2^x,3^x)$.  The global implication is a clean
application of the support problem, while the local equivalence is elementary finite-group
algebra.

The main quantitative advance is the Kummer defect theorem.  Multiplicative independence does
not merely create infinitely many bad primes; it creates a Frobenius set of explicit positive
density $1/\ell^2$.  A hypothetical counterexample therefore has positive-density local
failures on all three natural rays.  These failures inject into interpolation denominators
through a general collision-valuation inequality, yielding controlled squarefree divisors and
maximal secant-denominator escape.

The remaining problem is now stated in a more rigid form.  One must convert the real equality
of logarithmic slopes into either almost-everywhere finite-field synchronization or an
archimedean object too small to absorb the positive-density defect-prime divisor.  The present
results identify the exact local data such an object must preserve, the scale at which current
box interpolation detects it, and the quantitative improvements needed for a contradiction.

\appendix
\section{Exact reconnaissance source}
\label{app:code}

\begin{lstlisting}[style=pythonstyle]
#!/usr/bin/env python3
"""Exact reconnaissance for congruence-kernel defect primes.

All arithmetic decisions use Python integers and pow(..., mod). No floating-point
value enters a primality, residue, order, or synchronization decision.
"""
from __future__ import annotations

from math import gcd, isqrt, log


def primes_upto(n: int) -> list[int]:
    sieve = bytearray(b"\x01") * (n + 1)
    sieve[:2] = b"\x00\x00"
    for q in range(2, isqrt(n) + 1):
        if sieve[q]:
            sieve[q*q:n+1:q] = b"\x00" * (((n - q*q) // q) + 1)
    return [q for q in range(2, n + 1) if sieve[q]]


def factor(n: int) -> list[int]:
    ans: list[int] = []
    q = 2
    while q*q <= n:
        if n % q == 0:
            ans.append(q)
            while n % q == 0:
                n //= q
        q += 1 if q == 2 else 2
    if n > 1:
        ans.append(n)
    return ans


def order_mod(a: int, p: int) -> int:
    assert gcd(a, p) == 1
    d = p - 1
    for q in factor(p - 1):
        while d % q == 0 and pow(a, d // q, p) == 1:
            d //= q
    return d


def lcm(a: int, b: int) -> int:
    return a // gcd(a, b) * b


def synchronized_exponent(M: int, A: int, p: int) -> int | None:
    """Return e modulo |<2,3>| with 2^e=M and 3^e=A mod p, if one exists."""
    if (6 * M * A) % p == 0:
        return None
    h = lcm(order_mod(2, p), order_mod(3, p))
    for e in range(h):
        if pow(2, e, p) == M % p and pow(3, e, p) == A % p:
            return e
    return None


def first_ray_defect(B: int, C: int, primes: list[int]):
    for p in primes:
        if (B * C) % p == 0:
            continue
        d = order_mod(B, p)
        if pow(C, d, p) != 1:
            return p, d
    return None


def pair_table(limit_first: int = 100_000, limit_count: int = 10_000) -> None:
    pfirst = primes_upto(limit_first)
    pairs = [(4, 9), (8, 27), (4, 27), (8, 9), (5, 13), (10, 38)]
    print("PAIR TABLE")
    for M, A in pairs:
        d2 = first_ray_defect(2, M, pfirst)
        d3 = first_ray_defect(3, A, pfirst)
        d6 = first_ray_defect(6, M*A, pfirst)
        first_full = None
        sync = good = 0
        for p in pfirst:
            if (6*M*A) % p == 0:
                continue
            if first_full is None and synchronized_exponent(M, A, p) is None:
                first_full = p
            if p <= limit_count:
                good += 1
                if synchronized_exponent(M, A, p) is not None:
                    sync += 1
        print(M, A, d2, d3, d6, first_full, f"{sync}/{good}")


def kummer_table(B: int = 2, C: int = 5, ell: int = 3) -> None:
    cutoffs = [10**3, 10**4, 10**5, 10**6]
    primes = primes_upto(cutoffs[-1])
    print("KUMMER TABLE", B, C, ell)
    for X in cutoffs:
        good = count = 0
        theta_good = theta_s = 0.0
        for p in primes:
            if p > X:
                break
            if (B*C) % p == 0:
                continue
            good += 1
            theta_good += log(p)
            if (p % ell == 1 and
                pow(B, (p-1)//ell, p) == 1 and
                pow(C, (p-1)//ell, p) != 1):
                count += 1
                theta_s += log(p)
        print(X, count, good, count/good, theta_s/X, theta_s/theta_good)


if __name__ == "__main__":
    pair_table()
    kummer_table()
\end{lstlisting}

\section{Proof-dependency ledger}
\label{app:ledger}

\begin{longtable}{p{0.31\textwidth}p{0.18\textwidth}p{0.42\textwidth}}
\toprule
Statement & Status & Dependency \\
\midrule
\endhead
Local synchronized-exponent criterion & \statusE & Cyclicity of $\F_p^\times$ and quotient
factorization. \\
Three-ray global rigidity & \statusC & Three applications of the 1997 support theorem plus
unique factorization. \\
Infinite defect sets & \statusC & Contrapositive of the support theorem. \\
Kummer independence & \statusE{} / classical & Valuation minor; Capelli irreducibility for
$X^\ell-r$; degree coprimality with $\Q(\zeta_\ell)$. \\
Density $1/\ell^2$ & \statusC & Kummer Galois group and Chebotarev density. \\
Defect order inequality & \statusE & Cyclic-group calculation of $\ell$-adic order. \\
Collision-valuation lemma & \statusE & $x-y\mid P(x)-P(y)$ in $\Z[X]$. \\
Secant divisibility & \statusE & Collision-valuation lemma. \\
Kummer squarefree divisor & \statusC & Chebotarev prime number theorem plus secant
divisibility. \\
Box denominator detection & \statusE & Translation of signed exponent vectors to two grid
points. \\
Secant-denominator limit & \statusC & Bugeaud--Corvaja--Zannier gcd theorem. \\
Finite tables & \statusD & Exact Python integer arithmetic; source and output digests stated. \\
Research targets & \statusP & No theorem asserted. \\
\bottomrule
\end{longtable}

\begin{thebibliography}{99}

\bibitem{UnifiedReport}
\emph{The Alaoglu--Erd\H{o}s Integer-Exponent Conjecture: Machine-Checked Partial Results,
Exact Conditional Structure, Determinant Methods, and the Remaining Barrier}, unified report
in the ProveIt corpus, August 21, 2026; source supplied with the present task.

\bibitem{AlaogluErdos1944}
L. Alaoglu and P. Erd\H{o}s,
On highly composite and similar numbers,
\emph{Transactions of the American Mathematical Society} \textbf{56} (1944), 448--469.

\bibitem{CorralesSchoof1997}
C. Corrales-Rodrig\'a\~nez and R. Schoof,
The support problem and its elliptic analogue,
\emph{Journal of Number Theory} \textbf{64} (1997), 276--290,
\href{https://doi.org/10.1006/jnth.1997.2114}{doi:10.1006/jnth.1997.2114}.

\bibitem{BCZ2003}
Y. Bugeaud, P. Corvaja, and U. Zannier,
An upper bound for the G.C.D. of $a^n-1$ and $b^n-1$,
\emph{Mathematische Zeitschrift} \textbf{243} (2003), 79--84,
\href{https://doi.org/10.1007/s00209-002-0449-z}{doi:10.1007/s00209-002-0449-z}.

\bibitem{Perucca2009}
A. Perucca,
The multilinear support problem for products of abelian varieties and tori,
\emph{International Mathematics Research Notices} (2010),
arXiv:\href{https://arxiv.org/abs/0903.0817}{0903.0817}.

\bibitem{Sgobba2023}
P. Sgobba,
Divisibility conditions on the order of the reductions of algebraic numbers,
\emph{Mathematics of Computation} \textbf{92} (2023), 2281--2305,
arXiv:\href{https://arxiv.org/abs/2110.08911}{2110.08911}.

\bibitem{SerreChebotarev}
J.-P. Serre,
\emph{Lectures on $N_X(p)$},
A K Peters/CRC Press, 2012.

\bibitem{Gao2026}
S. Gao,
Counterexamples to the Jacobian conjecture in dimensions greater than two,
arXiv:\href{https://arxiv.org/abs/2608.00222}{2608.00222}, 2026.

\bibitem{MengYang2026}
G. Meng and L. Yang,
A five-variable counterexample to the Hessian conjecture, and the low-dimensional status of
the Jacobian and Hessian conjectures,
arXiv:\href{https://arxiv.org/abs/2607.22198}{2607.22198}, 2026.

\bibitem{HadamardOEIS}
OEIS Foundation,
Sequence A007299, historical notes and August 2026 update on the twelve Hadamard orders,
\url{https://oeis.org/A007299}, accessed August 22, 2026.

\bibitem{EpochHadamard}
Epoch AI,
Hadamard Matrix of Order 668, solution update,
\url{https://epoch.ai/frontiermath/open-problems/hadamard}, accessed August 22, 2026.

\bibitem{EpochRank}
Epoch AI,
Elliptic Curves over $\Q$ of Large Rank,
\url{https://epoch.ai/frontiermath/open-problems/elliptic-curve-rank}, accessed August 22, 2026.

\end{thebibliography}

\end{document}
